\subsection{The Natural Logarithm} 
Because we often want to consider exponential equations base $e$, and because of another way we could have defined logarithms using calculus, the logarithm base $e$ has a special name and notation.
\begin{definition}[Natural Logarithm]
The log base $e$ is called the \term{natural logarithm} and denoted $\ln$.
\end{definition}
\begin{example}
Find $\ln(e^3)$.
\begin{itemize}
\item $\ln(e^3)$ means \makeblank[1in]
\item So by definition $\ln(e^3)=\makeblanksmall$
\end{itemize}
\end{example}
The calculator has a button for the natural logarithm.
\begin{example}
Find $\ln(31)$ on the calculator. Round your answer to two decimal places.
\begin{itemize}
\item $\ln(31)\approx \makeblanksmall$
\end{itemize}
\end{example}
The calculator also has an option to calculate powers of $e$. On TI calculators, it is the \makeblanksmall function of the \makeblanksmall button.
\begin{example}
Find $e^3$ on the calculator. Round your answer to three decimal places. Check by taking the natural log.
\begin{itemize}
\item On the calculator, $e^3\approx\makeblanksmall$
\item To check, we take $\ln\left(\makeblanksmall\right)\approx\makeblanksmall\approx\makeblanksmall$
\end{itemize}
\end{example}